% 独立预览文件；不会被 paper_overleaf/main.tex 引用。
\documentclass[UTF8,a4paper,zihao=-4]{ctexart}
\usepackage{geometry}
\usepackage{array}
\usepackage{booktabs}
\usepackage{tabularx}
\usepackage{url}
\usepackage[hidelinks]{hyperref}
\geometry{left=2.5cm,right=2.5cm,top=2.5cm,bottom=2.5cm}
\pagestyle{empty}

\begin{document}

\begin{table}[htbp]
\centering
\footnotesize
\caption{附件数据的结构与主要用途}
\label{tab:global-data-overview}
\setlength{\tabcolsep}{5pt}
\renewcommand{\arraystretch}{1.28}
\begin{tabularx}{\textwidth}{@{}
    >{\raggedright\arraybackslash}p{3.45cm}
    >{\raggedright\arraybackslash}p{2.30cm}
    >{\raggedright\arraybackslash}p{3.05cm}
    >{\raggedright\arraybackslash}X@{}}
\toprule
\textbf{数据文件} &
\textbf{数据规模与粒度} &
\textbf{主要关联字段} &
\textbf{主要建模用途} \\
\midrule
{\scriptsize\nolinkurl{workload_trace.xlsx}}
& $50\,000$条任务
& \texttt{TaskID}、\texttt{TaskType}、\texttt{SourceRegion}
& 提供任务到达时刻、GPU需求、执行时长、来源区域、时延阈值及完成期限 \\
\addlinespace[2pt]

{\scriptsize\nolinkurl{GPU_information.xlsx}}
& $6$个区域
& \texttt{Region}
& 提供可调度GPU容量、IT功率上限、PUE及设施侧功率上限 \\
\addlinespace[2pt]

{\scriptsize\nolinkurl{network_latency.xlsx}}
& $36$个区域对
& \texttt{FromRegion}、\texttt{ToRegion}
& 形成区域间网络时延矩阵，判断跨区域执行的时延可行性 \\
\addlinespace[2pt]

{\scriptsize\nolinkurl{power_mapping.xlsx}}
& $3$类任务
& \texttt{TaskType}
& 将任务占用的等效GPU数量转换为AI IT功率 \\
\addlinespace[2pt]

{\scriptsize\nolinkurl{region_time_data.xlsx}}
& $6\times2407$条区域--小时记录
& \texttt{Region}、\texttt{Hour}
& 提供逐时电价、售电价格、碳强度、可用新能源、非AI负荷及基准运行状态 \\
\addlinespace[2pt]

{\scriptsize\nolinkurl{storage_information.xlsx}}
& $6$个区域
& \texttt{Region}
& 提供储能容量、SOC边界、初始SOC、充放电功率、效率及购售电边界 \\
\bottomrule
\end{tabularx}
\end{table}

\end{document}
